\section{梯度与方向导数}
\label{sec:02-Calculus/06-Multivariable-Calculus/03}

几乎所有训练循环里都躺着同一行更新式：
\[
\mathbf w\leftarrow\mathbf w-\eta\,\nabla L(\mathbf w).
\]
课上通常一句``梯度是上升最快的方向，所以往反方向走''就带过去了。这句话是个定理，不是定义，而且它的证明里藏着一个平时不写出来的选择——你怎样度量``一步''的大小。把这个选择换掉，最速方向就跟着换，训练里那些看似工程味十足的花招（逐坐标缩放、预条件、自然梯度）会在同一个式子里现形。

先把断言本身弄清楚。站在曲面 \(z=f(x,y)=x^2+2y^2\) 的点 \((1,1)\) 上，海拔 \(3\)。两个偏导是 \(f_x=2x=2\)、\(f_y=4y=4\)：往东走一小步升 \(2\)，往北走同样长的一步升 \(4\)。那么往东北走呢，是不是 \((2+4)/2=3\)？

\subsection{方向导数把函数压回一根数轴上}

设 \(\mathbf u\) 是单位向量。方向导数
\[
D_{\mathbf u}f(\mathbf a)
=\lim_{t\to0}\frac{f(\mathbf a+t\mathbf u)-f(\mathbf a)}{t}
\]
就是把 \(f\) 限制在从 \(\mathbf a\) 出发、沿 \(\mathbf u\) 的那根射线上，得到单变量函数 \(g(t)=f(\mathbf a+t\mathbf u)\)，再取 \(g'(0)\)。结构和单变量导数一模一样，只是数轴被摆到了 \(n\) 维空间里。

对上面的例子取 \(\mathbf u=(1,1)/\sqrt2\)，直接展开 \(f(1+t/\sqrt2,\,1+t/\sqrt2)=3+\tfrac{6t}{\sqrt2}+\tfrac32t^2\)，一阶系数给出
\[
D_{\mathbf u}f(1,1)=\frac{6}{\sqrt2}\approx 4.24 .
\]
既不是 \(3\)，也没停在两个偏导之间——它比大的那个偏导 \(4\) 还大。最陡的方向显然不在坐标轴上。

上一节的可微性此刻兑现了它的价值。若 \(f\) 在 \(\mathbf a\) 可微，
\[
f(\mathbf a+t\mathbf u)-f(\mathbf a)=\nabla f(\mathbf a)^{\trans}(t\mathbf u)+r(t\mathbf u),
\]
两边除以 \(t\) 让 \(t\to0\)，余项按定义消失，剩下
\[
D_{\mathbf u}f(\mathbf a)=\nabla f(\mathbf a)^{\trans}\mathbf u
=\ip{\nabla f(\mathbf a)}{\mathbf u}.
\]
无穷多个方向的变化率，被 \(n\) 个数一次性编码完毕。要注意这条等式是可微性的产物：上一节那个 \(xy/\sqrt{x^2+y^2}\) 的例子里，所有方向导数都存在，梯度也算得出来（等于零向量），这条等式却是错的。梯度能代表所有方向，靠的不是它的定义，是可微性这份担保。

\subsection{最速上升方向来自一条不等式}

在 \(\norm{\mathbf u}=1\) 的约束下最大化 \(\ip{\nabla f}{\mathbf u}\)。Cauchy--Schwarz 给出
\[
\ip{\nabla f}{\mathbf u}\le\norm{\nabla f}\,\norm{\mathbf u}=\norm{\nabla f},
\]
等号当且仅当 \(\mathbf u\) 与 \(\nabla f\) 同向。于是梯度方向是上升最快的方向，最大上升率恰为 \(\norm{\nabla f}\)；负梯度方向下降最快；与梯度正交的方向上一阶变化为零。

回到 \((1,1)\)：\(\nabla f=(2,4)\)，\(\norm{\nabla f}=2\sqrt5\approx4.47\)，最陡方向是 \((1,2)/\sqrt5\)。东北方向的 \(4.24\) 之所以略小，是因为它与最陡方向差了个夹角；一般地 \(D_{\mathbf u}f=\norm{\nabla f}\cos\alpha\)，变化率随偏离角按余弦衰减。沿正 \(x\) 方向只有 \(2\)，那是把 \(y\) 方向更陡的坡完全放过了。

\begin{center}
\begin{tikzpicture}[x=1.5cm,y=1.5cm,font=\small,>=stealth]
  \draw[black!25] (0,0) ellipse (1 and 0.707);
  \draw[black!35] (0,0) ellipse (1.732 and 1.225);
  \draw[black!25] (0,0) ellipse (2.236 and 1.581);
  \draw[->,black!45] (-2.7,0) -- (2.7,0) node[right] {\(x\)};
  \draw[->,black!45] (0,-1.9) -- (0,1.95) node[above] {\(y\)};
  \draw[black!60,dashed] (0.195,1.402) -- (1.805,0.598);
  \draw[->,red!70!black,very thick] (1,1) -- (1.40,1.80);
  \draw[->,blue!70!black,thick] (1,1) -- (1.636,1.636);
  \draw[->,black!70,thick] (1,1) -- (1.9,1);
  \fill (1,1) circle (2pt);
  \node[right,red!70!black] at (1.30,1.92) {\(\nabla f\)：\(4.47\)};
  \node[right,blue!70!black] at (1.66,1.55) {东北：\(4.24\)};
  \node[right,black!70] at (1.92,0.98) {正 \(x\)：\(2\)};
  \node[right,black!60] at (1.70,0.50) {切向：\(0\)};
  \node[black!45] at (-1.35,1.30) {\(f=5\)};
\end{tikzpicture}

\emph{三条等高线来自 \(f=x^2+2y^2\)。梯度垂直穿过 \((1,1)\) 处的等高线，其余方向的变化率按与它的夹角余弦打折；沿等高线切向恰好折到零。等高线越挤，\(\norm{\nabla f}\) 越大。}
\end{center}

图里梯度与等高线垂直不是巧合。设曲线 \(\mathbf r(t)\) 始终待在同一条水平集上，\(f(\mathbf r(t))=c\)，两边对 \(t\) 求导（用的是下一节的链式法则）得
\[
\ip{\nabla f(\mathbf r(t))}{\mathbf r'(t)}=0 .
\]
\(\mathbf r'(t)\) 是水平集的切向量，而梯度与所有这样的切向量正交，所以梯度落在法向上。这条结论此后会反复用到：隐式曲线 \(F(x,y)=c\) 在 \(\nabla F\ne\mathbf 0\) 处的法向量就是 \(\nabla F\)，切线可以写成 \(\ip{\nabla F(a,b)}{(x-a,\,y-b)}=0\)，不必先把 \(y\) 解成 \(x\) 的显函数，也不怕切线竖直。下一节讲约束优化时，整套 Lagrange 几何都建在这句话上。

\subsection{``最速''取决于你怎样量一步}

现在回到开头埋的那个伏笔。上面的推导里，\(\norm{\mathbf u}=1\) 用的是欧氏范数，而这个选择完全可以换。若用一个正定矩阵 \(M\) 定义内积 \(\ip{\mathbf a}{\mathbf b}_M=\mathbf a^{\trans}M\mathbf b\) 和相应的范数，同样的问题
\[
\max_{\norm{\mathbf u}_M=1}\;Df(\mathbf a)\,\mathbf u
\]
的解不再是 \(\nabla f\)，而是 \(M^{-1}\nabla f\)（再归一化）。

\begin{center}
\begin{tikzpicture}[x=1.6cm,y=1.6cm,font=\small,>=stealth]
  \draw[black!55] (0,0) circle (1);
  \draw[black!55,dashed] (0,0) ellipse (1.7 and 0.55);
  \draw[->,black!35] (-2.0,0) -- (2.2,0);
  \draw[->,black!35] (0,-1.0) -- (0,1.6);
  \draw[black!30] (-0.712,1.566) -- (1.712,0.166);
  \draw[black!30] (0.444,0.869) -- (2.522,-0.331);
  \draw[->,red!70!black,very thick] (0,0) -- (0.5,0.866);
  \draw[->,blue!70!black,very thick] (0,0) -- (1.483,0.269);
  \fill (0,0) circle (1.6pt);
  \node[left,red!70!black] at (0.42,0.95) {\(\nabla f\)};
  \node[below,blue!70!black] at (1.30,0.18) {\(M^{-1}\nabla f\)};
  \node[black!55] at (-0.78,-0.72) {欧氏单位球};
  \node[black!55] at (1.75,-0.72) {\(M\) 单位球};
\end{tikzpicture}

\emph{两条细线是同一个线性函数 \(\mathbf u\mapsto Df(\mathbf a)\mathbf u\) 的等值线。在圆形单位球上，走得最远的点在 \(\nabla f\) 方向；换成扁平的椭球，触到最外侧等值线的点被拽向椭球伸展的那一侧。函数没变，度量变了，最速方向就变了。}
\end{center}

真正不依赖度量的对象是导数本身——那个把位移送到函数变化量的线性泛函 \(Df(\mathbf a)\)。梯度是它借助某个内积找到的向量代表。行向量到列向量这一步，必须经过内积；换内积，代表就换。

这件事在训练里天天发生。参数向量的各个分量量纲常常差几个数量级：嵌入层的一个坐标和最后一层偏置的一个坐标，凭什么用同一把尺子相加求平方和？欧氏范数默认它们等价，于是学习率必须迁就最敏感的那个方向。逐坐标自适应方法（如 Adam）实际上是在用一个对角 \(M\) 重新定义单位球；\hyperref[sec:09-Convex-Optimization/06-Optimization-Algorithms/02]{``Newton 法用曲率校正方向''}用的是 \(M=H\)；而当参数索引的是一族概率分布时，\hyperref[sec:12-Real-Analysis-Support/09-Information-Geometry/02]{``自然梯度沿分布几何下降''}用 Fisher 信息矩阵作 \(M\)，量的是分布之间的距离而非参数之间的距离。它们的差别不在``用不用梯度''，而在``一步有多大''这个问题的答案。

\begin{remark}
一个具体的坑：若方向向量 \(\mathbf v\) 没有归一化，\(\ip{\nabla f}{\mathbf v}\) 混进了移动速度的信息。问``每单位距离变化多少''时必须先算 \(\mathbf v/\norm{\mathbf v}\)。同理，报告梯度范数时如果不说清用的是哪个范数、参数怎样缩放，这个数几乎不可比。
\end{remark}

\subsection{梯度是什么，取决于函数是什么}

梯度这个操作吃一个标量场，吐一个向量场：空间的每一点得到一个向量。它只描述场自身的空间变化率，至于这个向量在具体模型里代表什么，要由模型另行规定。温度分布 \(T\) 的梯度给的是每单位距离的升温量，热流大致沿 \(-\nabla T\)，但流量还要乘热导率；势能 \(V\) 的梯度给的是势能变化率，保守力是 \(-\nabla V\)。同一个数学对象，配不同的物理常数才成为不同的物理量。

到这里，一阶信息的第一层已经就位：可微性保证存在统一的线性模型，梯度把这个模型编码成一个向量，最速方向由内积决定。但这套语言只覆盖标量输出。神经网络里绝大多数中间层的输出是向量，复合层数动辄上百，需要的是把``线性模型''从向量升级成矩阵、并且说清复合时它们怎样相乘。下一节做这件事。
